\begin{table}[htbp]
\caption{Pairwise DeLong Tests for Overall AUC (Learner-based Split)}
\label{tab:delong_overall}
\centering
\scriptsize
\setlength{\tabcolsep}{3pt}
\begin{tabular}{lccc}
\toprule
\textbf{Dataset} & \textbf{IRT vs DKT} & \textbf{IRT vs SimpleKT} & \textbf{DKT vs SimpleKT} \\
\midrule
ASSISTments 2012 & $p < 0.001^{*}$ & $p < 0.001^{*}$ & $p < 0.001^{*}$ \\
Junyi Academy    & $p < 0.001^{*}$ & $p < 0.001^{*}$ & $p < 0.001^{*}$ \\
XES3G5M          & $p < 0.001^{*}$ & $p < 0.001^{*}$ & $p < 0.001^{*}$ \\
\bottomrule
\multicolumn{4}{p{\linewidth}}{\textit{Note:} Bonferroni-corrected significance threshold: $\alpha = 0.05/9 = 0.0056$. $^{*}$ indicates significant after correction. IRT comparisons serve as the classical baseline reference; IRT's learner-based AUC remains at 0.5000 because unseen learners do not have estimated ability parameters; however, its ACC reflects majority-class and item/concept difficulty effects rather than discriminative ranking ability. DeLong tests are computed using paired prediction-level outputs aligned by test instances within each dataset.} \\
\end{tabular}
\end{table}
